% Complete Section 8.6, original PDF 113-114 / printed 103-104.
% block: maxwell-equations-heading page:113 type:section
\subsection{\BT{Maxwell's equations}{Maxwell 方程}}\label{sec:maxwell-equations}

% block: maxwell-traditional-system page:113 type:ordinary-paragraph
\BT{In this section we show how to formulate Maxwell's equations as, essentially, a Hodge wave equation. The equations are, in traditional notation,}{本节说明, 如何将 Maxwell 方程在本质上表述为 Hodge 波动方程. 用传统记号, 这些方程为}
\begin{equation}\label{eq:maxwell-ampere}D_t-\operatorname{curl}H=-j,\end{equation}
\begin{equation}\label{eq:maxwell-faraday}B_t+\operatorname{curl}E=0,\end{equation}
\begin{equation}\label{eq:maxwell-magnetic-divergence}\operatorname{div}B=0,\end{equation}
\begin{equation}\label{eq:maxwell-electric-divergence}\operatorname{div}D=q.\end{equation}

% block: maxwell-constitutive-form-interpretation page:113 type:ordinary-paragraph
\BT{We view the electric field $E$ and the magnetizing field $H$ as (vector proxies for) 1-forms and the magnetic field $B$ and electric displacement $D$ as 2-forms. Thus, both the curl operators in the equations are the exterior derivative $d^1$ and both the divergence operators are $d^2$. The current density $j$ is thus a 2-form, and the charge density $q$ a 3-form. These equations are supplemented by the constitutive equations}{将电场 $E$ 和磁化场 $H$ 视为 1-形式 (的向量代理), 将磁场 $B$ 和电位移 $D$ 视为 2-形式. 因而, 方程中的两个旋度算子都是外微分 $d^1$, 两个散度算子都是 $d^2$. 电流密度 $j$ 因而是 2-形式, 电荷密度 $q$ 是 3-形式. 这些方程还需补充本构方程}
\[
D=\epsilon E,\qquad B=\mu H,
\]
\BT{where the permittivity $\epsilon$ and permeability $\mu$ are positive scalars or symmetric positive definite operators which may vary in space (but not time). Note that equation $D=\epsilon E$ has a clear meaning if we view both $D$ and $E$ as vector fields but less so if we think of it as relating differential forms of different orders. Written correctly in the language of differential forms, the equation is $D=\star_\epsilon E$, where the Hodge star operator $\star_\epsilon$ is formed with respect to an $\epsilon$-weighted inner product rather than the usual unweighted $L^2\Lambda^k$ inner product.}{其中介电常数 $\epsilon$ 和磁导率 $\mu$ 是正标量或对称正定算子, 可以随空间变化 (但不随时间变化). 注意, 若将 $D$ 和 $E$ 都视为向量场, 方程 $D=\epsilon E$ 的意义很清楚; 若将其视为联系不同次数微分形式的方程, 意义就不那么清楚. 用微分形式语言正确地写出, 该方程应为 $D=\star_\epsilon E$, 其中 Hodge 星算子 $\star_\epsilon$ 是关于 $\epsilon$-加权内积而非通常的无权 $L^2\Lambda^k$ 内积构成的.}

% block: maxwell-data-compatibility page:113 type:ordinary-paragraph
\BT{We note that \eqref{eq:maxwell-ampere} and \eqref{eq:maxwell-electric-divergence} imply a necessary compatibility condition on the data, namely, that $q_t=-\operatorname{div}j$. In the notation of exterior calculus, this says that $q_t=-dj$ or, weakly,}{注意, \eqref{eq:maxwell-ampere} 和 \eqref{eq:maxwell-electric-divergence} 蕴含数据的一个必要相容条件, 即 $q_t=-\operatorname{div}j$. 用外微积分记号, 就是 $q_t=-dj$, 或写成弱形式,}
\begin{equation}\label{eq:maxwell-data-compatibility}
\int q_t\wedge\tau-\int j\wedge d\tau=0,\qquad\tau\in\mathring H\Lambda^0.
\end{equation}

% block: maxwell-initial-boundary-data page:113 type:ordinary-paragraph
\BT{The initial data we impose on $E$, $D$, $B$, and $H$ must be compatible as well. Obviously we must have}{对 $E$, $D$, $B$ 和 $H$ 施加的初始数据也必须相容. 显然必须有}
\[
D(0)=\epsilon E(0),\qquad B(0)=\mu H(0),
\]
\BT{and, from \eqref{eq:maxwell-magnetic-divergence} and \eqref{eq:maxwell-electric-divergence}, we require}{并且由 \eqref{eq:maxwell-magnetic-divergence} 和 \eqref{eq:maxwell-electric-divergence}, 还要求}
\[
\operatorname{div}B(0)=0,\qquad\operatorname{div}D(0)=q(0),
\]
\BT{as well. Finally the Maxwell problem is completed with boundary conditions. For simplicity, we take homogeneous electric (perfect conductor) conditions $E\times n=0$, $B\cdot n=0$, i.e., $\operatorname{tr}E=0$ and $\operatorname{tr}B=0$.}{成立. 最后, 加上边界条件使 Maxwell 问题完整. 为简单起见, 取齐次电边界 (理想导体) 条件 $E\times n=0$, $B\cdot n=0$, 即 $\operatorname{tr}E=0$ 和 $\operatorname{tr}B=0$.}

% block: maxwell-slack-variable-and-weights pages:113,114 type:ordinary-paragraph
\BT{Our immediate goal is to write Maxwell's equations in the form of the 1-form Hodge wave equation. The unknown time-dependent 1-form, denoted by $v$ in \eqref{eq:hodge-wave-weak-formulation}, will be the electric field $E$. The unknown 2-form, denoted by $\beta$ above, will be the negative of the magnetic field $B$. The remaining dependent variable, the 0-form $\sigma$, will in fact turn out to vanish identically. It is a slack variable, whose vanishing reflects the compatibility of the data $q$ and $j$. (Our weak formulation will be well-posed even when the data $q$ and $j$ are incompatible, but, in that case, $\sigma$ will not vanish, and $E$ and $B$ will not satisfy Maxwell's equations.) As explained above, the variable coefficients $\epsilon$ and $\mu$ are handled by adjusting the inner product we use on the spaces of $L^2$ differential forms. For differential 1-forms we weight the inner product by $\epsilon$,}{当前目标是将 Maxwell 方程写成 1-形式 Hodge 波动方程的形式. 未知的含时 1-形式, 在 \eqref{eq:hodge-wave-weak-formulation} 中记为 $v$, 将是电场 $E$. 上面记为 $\beta$ 的未知 2-形式, 将是磁场 $B$ 的负值. 剩下的因变量, 即 0-形式 $\sigma$, 实际上会恒等于零. 它是松弛变量, 其为零反映了数据 $q$ 和 $j$ 的相容性. (即使数据 $q$ 和 $j$ 不相容, 我们的弱形式仍适定; 但此时 $\sigma$ 不为零, $E$ 和 $B$ 也不满足 Maxwell 方程.) 如上所述, 变系数 $\epsilon$ 和 $\mu$ 通过调整 $L^2$ 微分形式空间上使用的内积来处理. 对微分 1-形式, 用 $\epsilon$ 给内积加权,}
\[
\langle E,F\rangle_\epsilon:=\int\star_\epsilon E\wedge F,
\]
\BT{which is simply $\int\epsilon E\cdot F\,dx$ when written in terms of the vector proxies. Similarly for differential 2-forms we weight the inner product by $\mu^{-1}$. For the 0-forms we may weight with an arbitrary positive function $\alpha$.}{用向量代理写出, 它就是 $\int\epsilon E\cdot F\,dx$. 类似地, 对微分 2-形式, 用 $\mu^{-1}$ 给内积加权. 对 0-形式, 可以用任意正函数 $\alpha$ 加权.}

% block: maxwell-weak-system page:114 type:ordinary-paragraph
\BT{We are now ready to write the Maxwell system in weak form. With $\sigma\equiv0$, we get from \eqref{eq:maxwell-ampere}--\eqref{eq:maxwell-electric-divergence} that $(\sigma,E,B):[0,T]\to\mathring H\Lambda^0\times\mathring H\Lambda^1\times\mathring H\Lambda^2$ satisfy}{现在可以将 Maxwell 系统写成弱形式. 取 $\sigma\equiv0$, 由 \eqref{eq:maxwell-ampere}--\eqref{eq:maxwell-electric-divergence} 得到 $(\sigma,E,B):[0,T]\to\mathring H\Lambda^0\times\mathring H\Lambda^1\times\mathring H\Lambda^2$ 满足}
\begin{equation}\label{eq:maxwell-weak-slack}
\langle\sigma_t,\tau\rangle_\alpha-\langle E,d\tau\rangle_\epsilon=\int q\wedge\tau,\qquad\tau\in\mathring H\Lambda^0,\ t\in[0,T],
\end{equation}
\begin{equation}\label{eq:maxwell-weak-electric}
\langle E_t,F\rangle_\epsilon+\langle d\sigma,F\rangle_\epsilon-\langle B,dF\rangle_{\mu^{-1}}=-\int j\wedge F,\quad F\in\mathring H\Lambda^1,\ t\in[0,T],
\end{equation}
\begin{equation}\label{eq:maxwell-weak-magnetic}
\langle B_t,C\rangle_{\mu^{-1}}+\langle dE,C\rangle_{\mu^{-1}}=0,\qquad C\in\mathring H\Lambda^2,\ t\in[0,T],
\end{equation}
\BT{with the initial conditions $\sigma(0)=0$, $E(0)=E_0$, $B(0)=B_0$. After changing variables from $B$ to $-B$, this problem is exactly of the form of the Hodge wave equation \eqref{eq:hodge-wave-weak-formulation}, except that the inner products are weighted, and the right-hand side of the first equation need not vanish. The well-posedness of the problem can then be established in the same way as for \eqref{eq:hodge-wave-weak-formulation}.}{初始条件为 $\sigma(0)=0$, $E(0)=E_0$, $B(0)=B_0$. 将变量 $B$ 换为 $-B$ 后, 此问题就恰好具有 Hodge 波动方程 \eqref{eq:hodge-wave-weak-formulation} 的形式, 区别仅在于内积是加权的, 且第一个方程的右端不必为零. 因而可用与 \eqref{eq:hodge-wave-weak-formulation} 相同的方法建立问题的适定性.}

% block: maxwell-slack-variable-vanishing page:114 type:ordinary-paragraph
\BT{Let us show that, under the assumption of compatible data, the slack variable $\sigma$ vanishes and that, with $D=\star_\epsilon E$ and $H=\star_{\mu=1}B$, the quadruple $(E,B,H,D)$ satisfies Maxwell's equations. To see that $\sigma$ vanishes, we first note that it vanishes at time $t=0$ by assumption. Next, we evaluate \eqref{eq:maxwell-weak-slack} at time $t=0$ and use the fact that $\star_\epsilon E(0)=D(0)$ and $dD(0)=q(0)$ to conclude that $\sigma_t(0)=0$. Next, we differentiate \eqref{eq:maxwell-weak-slack} in time and add it to \eqref{eq:maxwell-weak-electric}, taking $F=d\tau$ in the latter. Using the compatibility condition \eqref{eq:maxwell-data-compatibility}, we obtain}{下面证明, 假定数据相容, 则松弛变量 $\sigma$ 为零, 并且取 $D=\star_\epsilon E$ 和 $H=\star_{\mu=1}B$ 后, 四元组 $(E,B,H,D)$ 满足 Maxwell 方程. 为说明 $\sigma$ 为零, 先注意由假设, 它在时刻 $t=0$ 为零. 接着, 在时刻 $t=0$ 考察 \eqref{eq:maxwell-weak-slack}, 并利用 $\star_\epsilon E(0)=D(0)$ 和 $dD(0)=q(0)$, 得出 $\sigma_t(0)=0$. 然后对 \eqref{eq:maxwell-weak-slack} 关于时间求导, 并与 \eqref{eq:maxwell-weak-electric} 相加, 在后一个方程中取 $F=d\tau$. 利用相容条件 \eqref{eq:maxwell-data-compatibility}, 得到}
\[
\langle\sigma_{tt},\tau\rangle_\alpha+\langle d\sigma,d\tau\rangle_\epsilon=0,\qquad\tau\in\mathring H\Lambda^0,\ t\in[0,T].
\]
\BT{Taking $\tau=\sigma_t$, we obtain that}{取 $\tau=\sigma_t$, 得到}
\[
\frac d{dt}(\|\sigma_t\|_\alpha^2+\|d\sigma\|_\epsilon^2)=0.
\]
\BT{Thus the quantity in parentheses is constant in time, and we have already seen that it vanishes when $t=0$. Therefore, $\sigma_t$ vanishes, i.e., $\sigma$ is constant in time, and so it too must vanish.}{因此括号内的量不随时间改变, 而我们已经看到它在 $t=0$ 时为零. 所以 $\sigma_t$ 为零, 即 $\sigma$ 不随时间改变, 因而它也必须为零.}

% block: maxwell-equations-recovered page:114 type:ordinary-paragraph
\BT{Having established that $\sigma$ vanishes, it is easy to see that \eqref{eq:maxwell-weak-slack}--\eqref{eq:maxwell-weak-magnetic} imply the Maxwell equations. Equation~\eqref{eq:maxwell-weak-slack} is a weak restatement of \eqref{eq:maxwell-electric-divergence}, \eqref{eq:maxwell-weak-electric} of \eqref{eq:maxwell-ampere}, and \eqref{eq:maxwell-weak-magnetic} of \eqref{eq:maxwell-faraday}. Finally \eqref{eq:maxwell-magnetic-divergence} follows from \eqref{eq:maxwell-faraday} since $\operatorname{div}B(0)=0$.}{已经证明 $\sigma$ 为零后, 容易看出 \eqref{eq:maxwell-weak-slack}--\eqref{eq:maxwell-weak-magnetic} 蕴含 Maxwell 方程. 方程 \eqref{eq:maxwell-weak-slack} 是 \eqref{eq:maxwell-electric-divergence} 的弱形式重述, \eqref{eq:maxwell-weak-electric} 是 \eqref{eq:maxwell-ampere} 的弱形式重述, 而 \eqref{eq:maxwell-weak-magnetic} 是 \eqref{eq:maxwell-faraday} 的弱形式重述. 最后, 由于 $\operatorname{div}B(0)=0$, \eqref{eq:maxwell-magnetic-divergence} 可由 \eqref{eq:maxwell-faraday} 推出.}
